\section{Conclusion}
\label{sec:conclusion}

This work presented a systematic study of knowledge conflicts in multi-hop QA across four LLMs, two datasets, three conflict types, and 19{,}416 examples. Five findings hold across models: (1)~answer-bearing hop conflicts reduce accuracy by 40--52pp on 2-hop and 44.0--46.5pp on 3-hop chains, all significant at $p < 0.0001$ (McNemar); (2)~answer-hop conflicts are far more damaging than bridge-hop ones (CFR $\leq$2.1\% at bridge, 27--36\% at answer); (3)~in 3-hop chains, final-hop conflicts cause near-total failure (Gemini 0.5\%) while three of four models show no significant early-hop drop (Gemini the sole exception at $\sim$13pp), revealing that models are disproportionately dependent on the final-answer document; (4)~within the Llama family, scale predicts post-conflict robustness (70B > 8B), but across families it does not (Qwen-32B and Gemini both underperform 8B Llama)---scale predicts the post-conflict \emph{floor}, not the drop size; (5)~temporal conflicts cause the largest absolute drop (50.8pp) and numerical match on relative drop (58.6\% vs.\ 57.0\%), while POR and baseline accuracy rise monotonically with Wikipedia pageview popularity in every model.

\paragraph{Implications.} For RAG system design, document quality control should prioritise the final-answer document; conflict detection should focus particularly on temporal and numerical claims (followed 50--80\% of the time under conflict).

\paragraph{Limitations.} Entity substitution produces surface-form conflicts only; the gold-context setting excludes HotpotQA's 8 distractor paragraphs and MuSiQue's additional candidates, so baselines are not directly comparable to published distractor-setting numbers (internal contrasts are unaffected). Only English bridge/compositional questions are evaluated; the popularity analysis excludes 160 of 866 factual entities with no Wikipedia article; API-based evaluation prevents mechanistic analysis.

\paragraph{Future Work.} The preliminary study sets up two follow-on phases. \emph{Semester~3 (conflict detection):} a BERT-based classifier trained on the 19{,}416 labelled examples from this study to predict whether a (question, retrieved-context) pair contains a conflict, using conflict-type (\S\ref{sec:types}) and popularity (\S\ref{sec:entity_popularity}) signals as features, evaluated by AUROC. \emph{Semester~4 (conflict mitigation):} contrastive decoding, self-consistency voting, hop-level backtracking, and selective abstention, fed by the detector and evaluated on the same benchmark with ablations. Further directions: extension to 4+ hop chains, simultaneous multi-hop injection, and mechanistic attention analysis.
